\section*{Problemas}

\subsection*{Enunciados de los problemas}

% Recordar
\begin{problema}
	\label{prob:287fdc5}
	Sea $X$ el número de ensayos Bernoulli independientes, con probabilidad de éxito $p$, requeridos hasta obtener el primer éxito (incluyendo el ensayo exitoso). Enuncie, sin derivar, la función de masa de probabilidad $f_X(k) = P(X=k)$ y su soporte $\mathcal{S}_X$.
\end{problema}

% Comprender
\begin{problema}
	\label{prob:0ebda90}
	Una transacción de base de datos falla de manera independiente en cada intento con probabilidad de fracaso $q=0.04$, y ya ha fallado en sus primeros $5$ intentos. Sin calcular ninguna probabilidad numérica, explique por qué la probabilidad de que necesite más de $8$ intentos en total para tener éxito, condicionada a haber fallado los primeros $5$, debe ser idéntica a la probabilidad incondicional de necesitar más de $3$ intentos desde el inicio del proceso.
\end{problema}

% Aplicar
\begin{problema}
	\label{prob:6ccfa13}
	Una firma de auditoría inspecciona transacciones de manera secuencial e independiente; una fracción $p=0.10$ presenta alguna irregularidad fiscal. Los auditores inspeccionan hasta encontrar exactamente $r=3$ transacciones irregulares. Calcule la probabilidad de que la tercera irregularidad se detecte exactamente en la vigésima inspección ($X=20$), y escriba la función de masa general de $X \sim \text{BN}(r,p)$ sobre su soporte.
\end{problema}

% Analizar
\begin{problema}
	\label{prob:a4a72f3}
	Sea $X \sim \text{BN}(r,p)$ el número de ensayos Bernoulli hasta acumular $r$ éxitos.
	\begin{enumerate}
		\item Demuestre por descomposición estocástica que $X = \sum_{i=1}^r Y_i$, donde $Y_i \sim \text{Geom}(p)$ son i.i.d.
		\item Derive la función generadora de momentos $M_{Y_1}(t) = \dfrac{pe^t}{1-qe^t}$ de una variable geométrica elemental.
		\item Usando independencia, deduzca $M_X(t)$ de la Binomial Negativa.
	\end{enumerate}
\end{problema}

% Evaluar
\begin{problema}
	\label{prob:ca5c4c8}
	Un ingeniero observa que el número de reintentos de conexión de un servidor tiene media muestral $\bar{x}=3.2$ y varianza muestral $s^2=13.44$, y propone modelarlo con una distribución de Poisson pura. Evalúe si esta elección es adecuada, calculando el índice de dispersión $s^2/\bar{x}$ y contrastándolo con el supuesto de equidispersión de Poisson ($\sigma^2=\mu$); de no serlo, justifique por qué un modelo Binomial Negativo (vía $Y=X-r$) es más apropiado.
\end{problema}

% Crear
\begin{problema}
	\label{prob:4c2c37d}
	Diseñe un ejemplo original (contexto, probabilidad de éxito $p$ y número de éxitos requeridos $r$) que se modele naturalmente con una distribución Geométrica ($r=1$) o Binomial Negativa ($r>1$), distinto de los ejemplos de bases de datos, auditoría o manufactura ya vistos en esta sección. Plantee la pregunta de probabilidad de interés y calcule $\mathbb{E}[X]$ y $\text{Var}(X)$ para los parámetros que eligió.
\end{problema}

\subsection*{Soluciones detalladas}

% Recordar
\begin{solproblema}
	Para que el primer éxito ocurra exactamente en el ensayo $k \ge 1$, los primeros $k-1$ deben ser fracasos y el ensayo $k$ un éxito: $f_X(k) = q^{k-1}p = (1-p)^{k-1}p$, para $k \in \mathcal{S}_X = \set{1,2,3,\dots}$.
\end{solproblema}

% Comprender
\begin{solproblema}
	La distribución geométrica tiene la propiedad de pérdida de memoria: $P(X > m+n \mid X>m) = P(X>n)$ para cualesquiera enteros $m,n\ge1$, ya que $P(X>k)=q^k$ depende únicamente de cuántos intentos adicionales se requieren, no del historial de fallos ya ocurridos. Con $m=5$ y $n=3$, se sigue directamente que $P(X>8\mid X>5) = q^{8-5} = q^3 = P(X>3)$: haber fallado ya $5$ veces no altera la probabilidad de necesitar $3$ o más intentos adicionales, porque el proceso "no recuerda" cuántos fracasos ya ocurrieron.
\end{solproblema}

% Aplicar
\begin{solproblema}
	Para que la tercera irregularidad ocurra exactamente en la inspección $20$, deben detectarse exactamente $r-1=2$ irregularidades en las primeras $19$ inspecciones (con probabilidad $\comb{19}{2}p^2q^{17}$) y la inspección $20$ debe ser irregular (probabilidad $p$):
	\begin{align*}
		P(X=20) = \comb{19}{2}p^3q^{17} = \comb{19}{2}(0.10)^3(0.90)^{17} \approx 0.0285.
	\end{align*}
	Generalizando, la función de masa de la Binomial Negativa es $f_X(k) = \comb{k-1}{r-1}p^r q^{k-r}$, para $k \in \mathcal{S}_X = \set{r, r+1, r+2, \dots}$.
\end{solproblema}

% Analizar
\begin{solproblema}
	\begin{enumerate}
		\item Sea $Y_1$ el número de ensayos hasta el primer éxito ($Y_1\sim\text{Geom}(p)$), y para $i\ge2$ sea $Y_i$ el número de ensayos adicionales entre el éxito $(i-1)$ y el éxito $i$. Por independencia de los ensayos Bernoulli y la propiedad de pérdida de memoria, cada $Y_i$ es independiente de los anteriores y sigue la misma distribución $\text{Geom}(p)$. El tiempo total hasta el $r$-ésimo éxito es la suma de estas $r$ esperas: $X = \sum_{i=1}^r Y_i$.
		\item Por definición, $M_{Y_1}(t) = \sum_{k=1}^\infty e^{tk}q^{k-1}p = pe^t\sum_{j=0}^\infty (qe^t)^j = \dfrac{pe^t}{1-qe^t}$, válida para $t < -\ln(q)$ (convergencia de la serie geométrica).
		\item Como $X$ es suma de $r$ variables i.i.d., su FGM es el producto de las FGM individuales: $M_X(t) = \left(\dfrac{pe^t}{1-qe^t}\right)^r$.
	\end{enumerate}
\end{solproblema}

% Evaluar
\begin{solproblema}
	El índice de dispersión es $s^2/\bar{x} = 13.44/3.2 = 4.2 \gg 1$: la varianza observada es más de $4$ veces la media, lo cual viola el supuesto de equidispersión de Poisson ($\sigma^2=\mu$). Modelar este tráfico como Poisson pura subestimaría sistemáticamente las colas de la distribución. Un modelo Binomial Negativo (vía $Y=X-r\sim\text{BN}(r,p)$, con $\mu_Y=r(1-p)/p$ y $\sigma_Y^2=r(1-p)/p^2$) es más apropiado porque tiene un parámetro adicional que absorbe explícitamente la sobredispersión: igualando momentos, $\sigma_Y^2/\mu_Y = 1/p$, de donde $\hat p = 3.2/13.44 \approx 0.2381$ y, sustituyendo en $\mu_Y$, $\hat r \approx 1$ — sugiriendo, de hecho, que el modelo Geométrico simple ya captura adecuadamente esta sobredispersión.
\end{solproblema}

% Crear
\begin{solproblema}
	\textbf{Ejemplo:} un equipo de ventas telefónicas cierra una llamada exitosamente con probabilidad $p=0.15$ por intento, de forma independiente. Sea $X$ el número de llamadas necesarias para lograr $r=4$ cierres exitosos, $X\sim\text{BN}(4, 0.15)$. Pregunta de interés: ¿cuántas llamadas se esperan hacer en promedio antes de lograr las $4$ ventas, y cuál es la variabilidad de ese conteo? $\mathbb{E}[X] = r/p = 4/0.15 \approx 26.67$ llamadas, y $\text{Var}(X) = r(1-p)/p^2 = 4(0.85)/(0.15)^2 \approx 151.11$, con desviación estándar $\sigma_X \approx 12.29$ llamadas.
\end{solproblema}
